\paragraph{My contribution.} I led the member intelligence part of the system. I defined the data-table contract that the other parts build against, and I built the data layer. It joins the Kaggle members with a synthetic scenario of clubs, classes, instructors, offers, past retention contacts and gym events, and it carries ground truth so that uplift and event effects can be measured. I built the feature engineering with its leakage check, the K-means segmentation (E1) and the calibrated churn model (E2). The model averages logistic regression and LightGBM, scores training members out of fold, and explains each score with per-member drivers. I also built the T-learner uplift baseline. On the app side I built the member and admin APIs, the data upload and rebuild flow, the dashboard and the member pages, including Member 360, which brings all three parts together. Finally, I reviewed the decision automation work, wrote the installation guide, and wrote the data, member intelligence and E1--E2 parts of this report.

\paragraph{What I learned.} The most useful lesson was that a good score is not the same as a good model. The churn model reached an F1 of 0.908 against 0.549 for a simple visit rule, but much of that came from a trend signal that compares this month's visits with the member's average. Without that signal and its close relatives, F1 fell to 0.817, which is the figure I would quote to a gym that measures visits differently. Calibration mattered as much as ranking. The planner multiplies probabilities by uplift and cost, so a well-ranked but badly calibrated score would have produced plans that look optimal and are not.

The scenario layer was the hardest part. The Kaggle data has no offers, no past contacts and no gym events, so I had to generate them, and everything generated is an assumption. I kept the ground truth separate from the training data, and I recorded which effects were invented and how strong they were. Even so, the uplift baseline recovered only part of the true ranking (Spearman 0.39 overall, 0.61 for free personal training), and the small contact history could not confirm whether it or the case-based reasoner gave better plans. Uplift is much harder to learn than risk, and an honest evaluation has to show that.

The fairness check was a lesson too. Recall for members aged 31 and over was 0.52 because they rarely leave, so the same model treats age groups differently without any change to its inputs. I now check groups separately as a matter of course.

\paragraph{Applying it.} From this project I take a habit of checking for leakage before trusting a strong result, calibrating any score that feeds a decision, and stating plainly what a synthetic evaluation can and cannot show. Agreeing on the data-table contract early also let the three parts work in parallel, and I would do that first on any team project.
